%!TEX root = ../template.tex
%%%%%%%%%%%%%%%%%%%%%%%%%%%%%%%%%%%%%%%%%%%%%%%%%%%%%%%%%%%%%%%%%%%%
%% Config/_package.tex
%% ISEL thesis configuration file
%%%%%%%%%%%%%%%%%%%%%%%%%%%%%%%%%%%%%%%%%%%%%%%%%%%%%%%%%%%%%%%%%%%%

\typeout{NT FILE Config/_packages.tex}%

% \nclangsetup{pt/contentsname=O MEU ÍNDICE}
% \nclangsetup{en/contentsname=MY LITTLE TOC}

\usepackage{pgfgantt}
\usepackage{xcolor}

% Color Definitions for the Gantt Chart
\definecolor{iselpurple}{RGB}{98, 43, 114} 
\definecolor{barblue}{RGB}{51, 102, 153}
\definecolor{groupblue}{RGB}{51, 102, 153}
\definecolor{linkred}{RGB}{165, 0, 33}
\definecolor{bargray}{RGB}{220, 220, 220}

%%============================================================
%%  FOR DEMO PURPOSES ONLY --- YOU SHOULD REMOVE THESE PACKAGES
%%============================================================
% \usepackage{float}        % FOR DEMO PURPOSES ONLY --- PLEASE REMOVE
% \restylefloat{figure}     % FOR DEMO PURPOSES ONLY --- PLEASE REMOVE


%%============================================================
%%  ADD YOUR PACKAGES AND OTHER CUSTOMIZATIONS HERE
%%============================================================

% \usepackage[orphans/prevent, widows/prevent]{widows-and-orphans} % add warnings about widows and orphans to “template.log”
% \usepackage{refcheck}     % list used and unused labels

\usepackage{makecell}      % Multiline table cells and control vertical and horizontal alignment

\usepackage{tcolorbox}    % pretty print color boxes

%------------------------------------------------------------------
% Customization of some packages
%------------------------------------------------------------------
\usepackage[T1]{fontenc}
\usepackage{lmodern}
%%% ---------------

%%% ------
% Where to look for figures
% \prependtographicspath{{some-folder-here}}

% Setup of listings, for more information check the 'listings' package manual
% \lstset{
%     captionpos=t,
%     basicstyle={\ttfamily\footnotesize},
%     numbers=left,
%     numberstyle={\ttfamily\tiny},
%     tabsize=2,
%     language=Java,
%     float,
%     frame=single,
% }

% --- Define Kotlin Language for Listings ---
\lstdefinelanguage{Kotlin}{
	keywords={package, as, typealias, class, interface, this, super, val, operator, var, fun, for, is, in, This, return, throw, when, it, import, while, if, else, null, try, break, continue, object, package, abstract, final, enum, open, annotation, override, private, public, protected, internal, implies, receiver, inc, impl, result, vararg, suspend, coroutine, data, companion},
	keywordstyle=\color{purple}\bfseries,
	ndkeywords={@Deprecated, @JvmName, @JvmStatic, @JvmOverloads, @JvmField, @JvmSynthetic, Iterable, Int, Integer, Float, Double, String, Runnable, dynamic},
	ndkeywordstyle=\color{orange}\bfseries,
	emph={println, return@, forEach, map, mapOf, listOf, orElse, binding, get, getReader, lines, reduce, trim, split, put, getOrDefault, add, getAsJsonObject, getAsString, split, write, readValue, service, setContentType, getWriter, deploy, build},
	emphstyle={\color{blue}},
	identifierstyle=\color{black},
	sensitive=true,
	commentstyle=\color{gray}\ttfamily,
	comment=[l]{//},
	morecomment=[s]{/*}{*/},
	stringstyle=\color{green!60!black}\ttfamily,
	morestring=[b]",
	morestring=[b]'
}

\lstdefinelanguage{yaml}{
	keywords={true,false,null,y,n},
	keywordstyle=\color{darkgray}\bfseries,
	basicstyle=\ttfamily\footnotesize, % basic font setting
	sensitive=false,
	comment=[l]{\#},
	morecomment=[s]{/*}{*/},
	commentstyle=\color{purple}\ttfamily,
	stringstyle=\color{blue}\ttfamily,
	morestring=[b]',
	morestring=[b]"
}

\lstdefinelanguage{json}{
	basicstyle=\ttfamily\footnotesize,
	numbers=left,
	numberstyle=\scriptsize,
	stepnumber=1,
	numbersep=8pt,
	showstringspaces=false,
	breaklines=true,
	frame=lines,
	stringstyle=\color{blue},
	keywordstyle=\color{black}\bfseries,
	morestring=[b]",
	literate=
	*{:}{{{\color{black}{:}}}}{1}
	{,}{{{\color{black}{,}}}}{1}
	{\{}{{{\color{black}{\{}}}}{1}
	{\}}{{{\color{black}{\}}}}}{1}
	{[}{{{\color{black}{[}}}}{1}
	{]}{{{\color{black}{]}}}}{1},
}


\usepackage{tikz}
\usetikzlibrary{shapes.geometric, arrows.meta, positioning, fit, calc, shadows}
\usepackage{graphicx}